% Material suplementario — Artículo 1 (español). Se compila antes que manuscript.tex.
\documentclass[11pt]{article}
\newcommand{\DocLang}{es}
\input{papers/shared/preamble.tex}
\input{papers/shared/authors.tex}
\hypersetup{pdftitle={Material suplementario: Seguridad poscomercializacion de VA-MENGOC-BC en Cuba}}

\begin{document}
\SupplementaryNumbering
\SyntheticNotice
\begin{center}
{\Large\bfseries Material suplementario\par}
\medskip
{\large Seguridad poscomercialización de la vacuna antimeningocócica B y~C VA-MENGOC-BC en Cuba,
\VStudyFirstYear--\VStudyLastYear: análisis nacional de notificaciones espontáneas con métodos de
desproporcionalidad, clases latentes y aprendizaje automático\par}
\medskip
{\small \AuthorOne{} y \AuthorTwo\par}
\end{center}

\tableofcontents

\section{Tasas anuales de notificación}\label{sec:S-rates}
\begin{table}[htbp]
\centering\small
\caption{Notificaciones anuales de ESAVI con VA-MENGOC-BC, dosis administradas y tasas de
notificación por \num{100000} dosis con intervalos de confianza de Poisson exactos del 95\,\%. La
razón compara cada tasa con la tasa esperada del programa, de \POExpectedRate{} por \num{100000}
dosis para todas las vacunas infantiles. Los años sin denominador de dosis muestran una raya.}
\label{tab:S-rates}
\POTableRates
\end{table}

\section{Selección del modelo de clases latentes}\label{sec:S-lca}
\begin{table}[htbp]
\centering\small
\caption{Modelos de clases latentes con $K=\POLCAKMin,\dots,\POLCAKMax$ clases ajustados a los
indicadores de eventos con prevalencia de al menos el 0,5\,\%. Se seleccionó el modelo con menor BIC;
el ICL-BIC penaliza las clases poco separadas y la $R^2$ de entropía mide la certeza de la
clasificación (1 = perfecta).}
\label{tab:S-lca}
\POTableLCASelection
\end{table}

\section{Factores asociados a la hospitalización}\label{sec:S-logistic}
\begin{table}[htbp]
\centering\small
\caption{Razones de odds ajustadas de hospitalización del modelo de regresión logística de Firth
multivariable ajustado a todas las notificaciones (intervalos de confianza de Wald del 95\,\%). Los
contrastes por región se refieren a la región occidental; los términos de eventos y vacunas son
indicadores de presencia frente a ausencia. Se suprimen los términos registrados en menos de
\VMinCell{} notificaciones.}
\label{tab:S-logistic}
\POTableLogistic
\end{table}

\section{Desproporcionalidad en cada diseño de comparación}\label{sec:S-designs}
Todas las tablas siguen la notación del texto principal: $a$, notificaciones con VA-MENGOC-BC y el
evento; $E$, valor esperado de $a$ bajo independencia; IC\textsubscript{025}, límite inferior de
credibilidad del 95\,\% del componente de información; EB\textsubscript{05}, percentil 5 inferior de
la media geométrica bayesiana empírica (solo en el diseño principal, porque la distribución previa
gamma-Poisson se estima con toda la base); Métodos, número de los cuatro métodos que cumplen sus
umbrales; Señal, criterio preespecificado (IC\textsubscript{025}${}>0$ con
$a\geq\PODispMinReports$). Los recuentos menores que \VMinCell{} se suprimen junto con los
estadísticos derivados de ellos.

\begin{table}[htbp]
\centering\small
\caption{Diseño principal: notificaciones con VA-MENGOC-BC frente al resto (todos los tipos de
evento).}
\label{tab:S-disp-primary}
\POTableDispPrimaryAllOtherVaccines
\end{table}

\begin{table}[htbp]
\centering\small
\caption{Diseño de comparador activo en lactantes: notificaciones en menores de \POInfantMonths{}
meses.}
\label{tab:S-disp-infant}
\POTableDispInfantActiveComparator
\end{table}

\begin{table}[htbp]
\centering\small
\caption{Diseño de una sola vacuna: VA-MENGOC-BC administrada sola frente a otras vacunas
administradas solas.}
\label{tab:S-disp-single}
\POTableDispExcludingCoadministration
\end{table}

\begin{table}[htbp]
\centering\small
\caption{Diseño de enmascaramiento: grupo de comparación sin las notificaciones que mencionan la
vacuna pentavalente (PENTA-L).}
\label{tab:S-disp-masking}
\POTableDispExcludingPentavalentMasking
\end{table}

\section{Estabilidad temporal del componente de información}\label{sec:S-cumic}
\begin{table}[htbp]
\centering\small
\caption{Componente de información (IC) acumulado con intervalo de credibilidad del 95\,\%,
recalculado con las notificaciones acumuladas hasta el final de cada año, para los cuatro eventos
más notificados y para cada evento que cumple el criterio principal de señal.}
\label{tab:S-cumic}
\POTableCumulativeIC
\end{table}

\section{Calidad de los datos}\label{sec:S-dq}
La calidad de los datos se evaluó con el marco armonizado de Kahn et al.\ (conformidad, completitud
y plausibilidad) sobre las notificaciones armonizadas antes de eliminar duplicados.

\begin{table}[htbp]
\centering\small
\caption{Completitud: porcentaje de notificaciones con cada campo clave registrado, por año del
fichero de origen.}
\label{tab:S-completeness}
\POTableCompleteness
\end{table}

\begin{table}[htbp]
\centering\small
\caption{Verificaciones de plausibilidad en todos los años. Cuando algún recuento anual se suprimió,
el número de registros señalados se muestra como cota inferior y se omite el porcentaje.}
\label{tab:S-plausibility}
\POTablePlausibility
\end{table}

\section{Enfermedad meningocócica en Cuba, \POITSFirstYear--\POITSLastYear}\label{sec:S-its}
\begin{figure}[htbp]
\centering
\includegraphics{p1_incidence}
\caption{Casos nacionales anuales de enfermedad meningocócica (escala logarítmica) y ajuste del
modelo binomial negativo segmentado. La banda sombreada marca la transición de \POITSTransition{}
(ensayo de eficacia, vacunación masiva e incorporación al programa nacional).}
\label{fig:S-its}
\end{figure}

\end{document}
